\documentclass[11pt]{article}
% preamble.tex — shared preamble for the Flyxion connective-essay series
% Compile target: XeLaTeX. \input this file, or copy its contents, at the
% top of each essay after \documentclass.

\usepackage{fontspec}
\setmainfont{TeX Gyre Termes} % Times-like serif; falls back gracefully if absent
\usepackage{amsmath,amssymb,amsthm}
\usepackage[margin=1.25in]{geometry}
\usepackage[hidelinks]{hyperref}
\usepackage{microtype}
\usepackage{enumitem}
\usepackage{footmisc}

\newtheorem{claim}{Claim}
\newtheorem{definition}{Definition}
\newtheorem{observation}{Observation}

% --- Corpus vocabulary macros -------------------------------------------
\newcommand{\admissible}{\textsc{admissible}}
\newcommand{\inadmissible}{\textsc{inadmissible}}
\newcommand{\repair}{\textit{repair}}
\newcommand{\continuation}{\textit{continuation}}
\newcommand{\rsvp}{\textsc{rsvp}}

% --- Title block helper ---------------------------------------------------
% Usage: \essaytitle{Title}{Subtitle or tagline}
\newcommand{\essaytitle}[2]{%
  \title{#1\\[0.3em]\large #2}
  \author{Flyxion\\\normalsize Independent Researcher}
  \date{\today}
}

\pagestyle{plain}

\essaytitle{Contract Law as an Admissibility Regime}{Formation, breach, and remedy read through continuation rather than snapshot validity}
\begin{document}
\maketitle
\begin{abstract}
Classical contract doctrine is usually taught in two stages. First comes formation: offer, acceptance, and consideration establish whether a contract exists. Only afterward come breach and remedy. That sequence is pedagogically convenient, but it obscures the deeper structure of the field. The better reading is that formation rules are an admissibility gate on the creation of binding obligation; material breach marks an inadmissible attempted continuation of that obligation; and remedies are repair operators that specify how the legal relationship can still go on once performance has failed. Expectation damages, in particular, are not best understood as merely downstream bookkeeping. They preserve the future the contract licensed. By contrast, reliance, restitution, and rescission preserve different invariants, and some doctrinal zones --- especially rescission for fraud or mutual mistake --- expose the limits of the repair framing by aiming to unwind the transaction as though the original transition should never have counted. Read this way, contract law sits naturally beside the admissibility-centered picture developed in \textit{Contradiction and Continuation}, while still retaining the doctrinal details that make contract law a legal practice rather than a metaphor.
\end{abstract}

\section{Introduction}
The black-letter common law of contract begins with formation. The Restatement (Second) of Contracts defines a contract as a promise or set of promises for the breach of which the law gives a remedy, or whose performance the law recognizes as a duty. It then states, in \S\S 17, 22, 24, and 71, that formation ordinarily requires a bargain, mutual assent, an offer and acceptance structure, and consideration.\footnote{\textit{Restatement (Second) of Contracts} \S\S 1, 17, 22, 24, 71 (Am. L. Inst. 1981).} In the usual presentation, those elements answer the primary question: is there a valid contract? Breach and remedy enter only once that question has already been settled.

That ordering is not false. A court often must decide whether any enforceable obligation formed before it can decide the legal consequences of nonperformance. But the ordering can be misleading if treated as conceptually basic. Contract law is not really organized around a static yes-or-no object called ``validity,'' with all later doctrine merely appended to it. It is organized around a regulated transition from no duty to duty, then around the legal conditions under which that relation can continue despite deviation, and finally around the modes of repair available when continuation in kind fails.

\subsection{Step 2: The unexamined premise}
The unexamined premise is that formation doctrine is the load-bearing center of the field, while breach and remedy are merely secondary consequences. On that premise, offer, acceptance, and consideration create the real legal object; everything after that is postscript.

The direction of explanation is closer to the reverse. Formation elements matter because they mark when the law will admit a change in normative position --- when the parties may be treated as having entered a binding, continuable relation rather than having merely talked, negotiated, boasted, or made a gratuitous promise. Breach doctrine and remedial doctrine are not detachable add-ons. They reveal what sort of thing a contract \emph{is}: a legally supervised trajectory of expected performance whose identity is partly constituted by what departures it tolerates and what repair mechanisms become available when it is not fully carried through.

\subsection{Step 3: Temporal versus constitutive priority}
Again the distinction is essential. Temporally, formation questions often come first in litigation and legal analysis. A court cannot award expectation damages on a bargain that never formed. In that sense, existence precedes breach.

Constitutively, however, the content of ``there is a contract here'' is inseparable from the law's account of continuation under stress. To say a contract exists is not merely to say that assent once occurred. It is to say that the legal system recognizes a structured obligation whose future deviations are assessable as substantial performance, immaterial breach, material breach, anticipatory repudiation, cure, waiver, damages, restitution, or rescission. In that sense the contract is not fully intelligible apart from the doctrines that govern its possible failure. The same point is central to the first part of \textit{Contradiction and Continuation}\footnote{Flyxion, \emph{Contradiction and Continuation} (draft), a study of Graham Priest's non-classical logics reworked as regimes of one constrained-representation theory, source at \url{https://github.com/standardgalactic/alphabet/blob/core/soup/book-latex-source/main.tex}, compiled PDF at \url{https://github.com/standardgalactic/alphabet/blob/core/soup/book-latex-source/contradiction-and-continuation.pdf}. Part I develops $x\rightsquigarrow_R y$ as consequence-as-preservation-under-admissible-transformation, cited here only for that narrow claim.}: a licensed transition is defined not only by how it starts but by what counts as preserving it under an admissible rule of transformation.

\section{Minimum apparatus: admissibility and repair}
Only two constructs are needed.

\begin{definition}
Let a contract state at time $t$ be the tuple $s_t=(f_t,h_t)$, where $f_t$ records formation-status and $h_t$ records the performance history up to $t$. Write $s_t \rightsquigarrow_{R_{\mathrm{contract}}} s_{t+1}$ when continuation from $t$ to $t+1$ is admitted by the contract rule $R_{\mathrm{contract}}$.\footnote{This is a direct specialization of Part I's general schema. \emph{Contradiction and Continuation}, Part I, defines consequence itself as $x \rightsquigarrow_R y$: $y$ is a licensed continuation of $x$ exactly when the transformation from $x$ to $y$ is admitted by a constraint structure $R$, not merely when $y$ follows $x$ in sequence. $R_{\mathrm{contract}}$ below is that same relation $R$, instantiated for one legal domain rather than borrowed as a loose metaphor.} That rule admits continuation iff (i) the purported agreement satisfies the formation conditions stated in \S\S 17, 22, 24, and 71 of the Restatement (Second) of Contracts, and (ii) at time $t$ no uncured material breach exists.\footnote{\textit{Restatement (Second) of Contracts} \S\S 17, 22, 24, 71; \textit{Restatement (Second) of Contracts} \S 241 (Am. L. Inst. 1981).} If either condition fails, ordinary contractual continuation is not admitted and some further doctrinal operator must specify what happens next.
\end{definition}

A \repair{} is then a doctrinally authorized continuation after failure of performance. Sometimes the law repairs by projecting forward to the promised future, as with expectation damages. Sometimes it repairs by preserving only reliance expenditures or stripping unjust enrichment, as with reliance or restitution. Sometimes it refuses continuation and tries to unwind the transaction altogether. Those are not merely different dollar amounts. They preserve different invariants.

A brief proof sketch shows that this is not just imported vocabulary. If the Restatement's formation conditions are not met, doctrine says no contract formed; in the notation above, $f_t$ never enters the formed state, so no $\rightsquigarrow_{R_{\mathrm{contract}}}$ continuation is admitted. If those conditions are met and no uncured material breach exists, the law treats the parties' duties as ongoing, which is exactly the claim that $s_t \rightsquigarrow_{R_{\mathrm{contract}}} s_{t+1}$. And once a material breach occurs, ordinary continuation fails unless cure, waiver, damages, or discharge re-specify the path. So $R_{\mathrm{contract}}$ recovers the field's own answers to ``did obligation form?'' and ``can this performance history still count as the same contract going on?'' rather than merely relabeling them from outside.

\section{Formation as admissibility: offer, acceptance, consideration}
The Restatement's formation provisions look much clearer once read as gatekeeping rules rather than as metaphysical ingredients of a mysterious legal object. An offer, under \S 24, is a manifestation of willingness to enter a bargain that justifies the other party in understanding that assent is invited and will conclude it. Acceptance, under \S 22, is ordinarily the complementary manifestation. Consideration, under \S 71, marks the bargain relation: each side's promise or performance is sought in exchange for the other.

Taken together, those rules decide when legal continuation may begin. Before they are satisfied, the law usually sees only negotiation or gratuitous intention. After they are satisfied, the parties occupy a new normative position in which future conduct will be assessed against a shared, public standard. That is why consideration is so important in the classical account. It is not just evidence that people were serious. It is a way of distinguishing a relation the law is prepared to carry forward from one it will leave in the moral or social domain.

This reading also clarifies why some formation doctrines are context-sensitive rather than absolute. Commercial settings, consumer settings, option contracts, promissory estoppel, and firm offers all adjust the admissibility gate. They do not merely generate exceptions to a pristine essence. They show that the law is always specifying which transitions into obligation it can stabilize under the pressures of actual practice.

\section{Mirror image and UCC \S 2-207: renegotiating the gate}
The common-law mirror-image rule is an excellent example. Under the classical rule, an acceptance had to mirror the offer. A purported acceptance that varied the terms was treated as a counteroffer, not as assent. This is doctrinally elegant because it keeps the admissibility gate crisp: no binding continuation unless the parties line up term for term.

But commercial practice exposed the cost of that elegance. Merchants routinely exchanged purchase orders, acknowledgments, invoices, and boilerplate forms whose printed terms did not match, even while everyone behaved as though a deal existed. To insist on perfect mirror imaging in such settings was to deny legal continuation to relationships that commerce had already operationally stabilized.

UCC \S 2-207 is therefore not a technical patch on the edges of formation doctrine. It is a renegotiation of the admissibility criterion itself. Section 2-207(1) provides that a definite and seasonable expression of acceptance can operate as an acceptance even though it states additional or different terms, unless acceptance is expressly made conditional on assent to those terms. Section 2-207(2) then treats additional terms between merchants as proposals that may become part of the contract unless the offer limited acceptance, the new term materially alters the bargain, or timely objection is given. Section 2-207(3) goes further still: conduct recognizing the existence of a contract can itself establish one, even when the writings do not.\footnote{U.C.C. \S 2-207.}

This is a remarkable doctrinal move. The law stops asking for exact textual symmetry and starts asking whether there is enough coordinated commercial conduct to license obligation anyway. The field did not discover that formation no longer mattered. It discovered that its old gate was too brittle for the transactions it claimed to govern. Put differently: a regime that insists on purity at the boundary can become \inadmissible{} to the very practice it regulates.

The battle of the forms also helps answer an obvious objection to the present essay. If formation is only an admissibility gate, does that trivialize assent? No. It shows why assent has to be made operational. The law must decide what sort of mismatch is small enough to preserve the identity of the transaction and what sort is serious enough to prevent or reconfigure continuation. UCC \S 2-207 is a live legal answer to that problem.

\section{Material breach and substantial performance as thresholds of continuation}
Once obligation exists, the next question is not simply whether there has been deviation. Deviation is ubiquitous. The deeper question is whether the deviation is still the same performance, legally speaking, or whether it crosses the threshold into a breach serious enough to suspend or discharge the other party's remaining duties.

That is exactly what the doctrine of substantial performance does. In \textit{Jacob \& Youngs, Inc. v. Kent} (N.Y. 1921), the builder installed pipe of the wrong brand but of comparable quality.\footnote{\textit{Jacob \& Youngs, Inc. v. Kent}, 230 N.Y. 239 (1921).} Cardozo refused to treat the deviation as grounds for tearing out the work and rebuilding the house to literal specification. The contract had, in the relevant sense, substantially been performed. The remedy tracked the difference in value rather than the cost of replacement.

This is a direct instance of admissibility reasoning. The law asks whether the deviation is compatible with treating the promised exchange as still going on under correction, or whether the deviation is so central that continuation in the same register is no longer licit. Restatement \S 241's familiar material-breach factors make the threshold explicit: how much expected benefit has been lost, whether that loss can be compensated, the risk of forfeiture, the likelihood of cure, and the role of good faith.\footnote{\textit{Restatement (Second) of Contracts} \S 241 (Am. L. Inst. 1981).}

Notice what substantial performance excludes. It excludes the fantasy that identity of performance requires byte-for-byte restoration of every contractual detail. A contract can continue through nontrivial variation. The law preserves the relation by attaching a compensatory adjustment rather than by declaring the whole performance a nullity. In the language of \textit{Distinction Holonomy}, this is why repair is not restoration. The path matters; what counts is whether the relation can be carried forward without importing a contradiction into the bargain's public meaning.

Material breach marks the opposite case. A material breach is not merely a wrong act added onto an otherwise intact relation. It is an \inadmissible{} transition in the attempted performance sequence: one party has moved in a way the existing obligation cannot license without triggering a new legal response, precisely because the deviation works an infliction of loss the bargain did not allocate to the non-breaching party. The other party may then suspend performance, treat the contract as breached, or seek a remedy calibrated to the failure.

Anticipatory repudiation sharpens the temporal point. In \textit{Hochster v. De La Tour}, the promisee was permitted to treat an announced refusal to perform as a present breach before the date of performance arrived.\footnote{\textit{Hochster v. De La Tour}, 2 El. \& Bl. 678 (1853).} The admissibility question is therefore assessed prospectively over a trajectory, not merely retrospectively over a completed act. Once one party repudiates in advance, the law judges that the future path no longer satisfies $R_{\mathrm{contract}}$: the threatened continuation is already \inadmissible{} as a way for this contract to go on.

\section{Remedies as distinct repair operators}
The most illuminating place to test the thesis is remedies. If remedies were merely secondary consequences, their internal differences would be economically important but conceptually incidental. In fact, they tell us what feature of the contractual relation the law is trying to preserve.

Expectation damages are the clearest example. Restatement \S 347 measures them by loss in value plus other foreseeable loss, minus costs avoided.\footnote{\textit{Restatement (Second) of Contracts} \S 347 (Am. L. Inst. 1981).} The standard textbook formula says they put the injured party in as good a position as if the contract had been performed. That is already a theory of repair. It projects the legal relation forward to the promised future and compensates so that the expected continuation is preserved as nearly as money can do it.

\textit{Hawkins v. McGee} (N.H. 1929), the ``hairy hand'' case, is famous precisely because it makes the point starkly.\footnote{\textit{Hawkins v. McGee}, 84 N.H. 114 (1929).} The measure is not simply what the plaintiff spent or lost in reliance. It is the difference between the value of the promised successful outcome and the value of what was actually received, adjusted by the ordinary limits on proof and causation. Expectation protects the contract's licensed future.

Reliance damages preserve a different invariant. They aim to restore the plaintiff to the position occupied before reasonable reliance on the promise. They are often especially salient where expectation is too speculative or where the legal wrong is promise-induced action rather than the lost benefit of the bargain. \textit{Sullivan v. O'Connor} and \textit{Chicago Coliseum Club v. Dempsey} each, in different ways, display the law's anxiety about speculative futures and its willingness to reimburse concrete preparation instead.\footnote{\textit{Sullivan v. O'Connor}, 363 Mass. 579 (1973); \textit{Chicago Coliseum Club v. Dempsey}, 265 Ill. App. 542 (1932).}

Restitution preserves yet another invariant: it strips benefits unjustly retained by the breaching party. It is less concerned with the plaintiff's promised future than with the defendant's gain. One can call this a repair only in a thinner sense. The legal system is not recreating the bargain's forward trajectory; it is preventing one party from keeping value for which the reciprocal structure failed.

Specific performance, where available, makes the structure even clearer. It is a repair operator that tries to continue the very performance itself rather than replacing it with money. Its exceptional status confirms the point rather than undermining it: some bargains are treated as inadequately repairable by ordinary damages because the continuation promised was too distinctive to cash out cleanly.

A serious rival account denies that expectation damages are a \repair{} operator in this sense at all. Holmes's ``bad man'' perspective and Posner's law-and-economics defense of efficient breach instead suggest that damages are the price of a lawful option: one may either perform or breach and pay, as though compensation were substitute performance rather than a response to failed continuation.\footnote{Oliver Wendell Holmes Jr., ``The Path of the Law,'' \emph{Harvard Law Review} 10, no. 8 (1897); Richard A. Posner, \emph{Economic Analysis of Law}, 9th ed. (New York: Wolters Kluwer, 2014), ch. 4 (efficient breach).} That view captures one real feature of a damages-centered system, namely that the law often prefers compensation to coercive supervision. But it overstates the optionality of breach. Material breach still counts as wrongful nonperformance; anticipatory repudiation lets the other party treat future nonperformance as a present failure; and specific performance remains available where money cannot preserve the right invariant. So damages are better understood as a repair operator the regime commonly selects after an inadmissible transition, not as proof that performance was optional all along.

\section{A companion point to \textit{Contradiction and Continuation}}
Read beside \textit{Contradiction and Continuation}, contract doctrine looks less like a field obsessed with static validity and more like a field for regulating permissible transformation. Part I's basic thought --- that consequence is licensed under an admissibility regime rather than delivered by bare sequence alone --- fits contract formation unusually well. A bargain does not become binding because one event merely follows another in time. It becomes binding when the law can treat the transition from negotiation to obligation as normatively coherent and publicly maintainable.

The deeper parallel appears at breach. A materially defective performance is not just a later fact appended to an already complete contract. It changes what continuations remain licit. Cure, waiver, damages, specific performance, restitution, and discharge are all ways the legal system re-specifies the path once the straightforward one has failed.

\section{Disconfirming instance: rescission really can unwind rather than repair}
The present framework has real limits, and contract law itself supplies them. Rescission for fraud or mutual mistake is often better described not as repair of a valid continuation but as refusal of the original transition. In such cases the law aims, so far as possible, to unwind the transaction and restore the parties to their pre-contract positions.

\textit{Sherwood v. Walker} (Mich. 1887), the famous ``barren cow'' case, is a classic illustration of mutual mistake undoing the deal's foundation.\footnote{\textit{Sherwood v. Walker}, 66 Mich. 568 (1887).} Whether one agrees with every aspect of the case, the doctrinal posture is clear: the law may treat the apparent contract as resting on such a mistaken basis that preservation of the promised future is the wrong aim. Likewise, fraud-based rescission does not seek to save the bargain's continuation. It seeks to deny that the law should carry forward a relation produced by a defective entry condition. In the formal terms introduced above, these are not merely cases where no $s_t \rightsquigarrow_{R_{\mathrm{contract}}} s_{t+1}$ can be found after breach. They are cases where $R_{\mathrm{contract}}$ is undefined on the purported input itself, because fraud or mutual mistake defeats the admissibility of the transition into obligation rather than merely interrupting a valid continuation.

A second, sharper qualification is comparative. The common law's tendency to treat money damages as the ordinary repair and specific performance as exceptional is not universal. Civil-law-oriented regimes often make performance primary, and CISG art. 46 gives the buyer a right to require performance that is much closer to a performance-first default than the common-law baseline.\footnote{United Nations Convention on Contracts for the International Sale of Goods (CISG), art. 46 (1980).} That means the repair taxonomy is itself regime-relative: what counts as the right continuation-preserving operator depends partly on which legal admissibility regime one is in.

This is a genuine counterexample to any overextended repair thesis. Sometimes the right description is not ``the relation persists under a new operator'' but ``the transition should never have been admitted on these terms.'' And sometimes the very ranking of repair operators shifts across legal regimes. Those distinctions keep the present argument honest. Not all contract doctrine is about maintaining the same relation through disturbance. Some of it is about declaring that the relation, properly understood, never earned continuation in the first place.

\section{Non-triviality conditions}
The essay's central claim would be empty if any of three things were true.

First, it would be empty if contract doctrine after formation were just remedial arithmetic appended to a fully self-sufficient object called a valid contract. But material breach, substantial performance, cure, waiver, and remedy selection all help determine what the contractual relation counts as under stress.

Second, it would be empty if the admissibility language merely redescribed obvious sequences without explanatory gain. The gain here is concrete: it clarifies why UCC \S 2-207 changes formation rules, why substantial performance tolerates deviation, and why expectation, reliance, restitution, and rescission are not interchangeable.

Third, it would fail if counterexamples collapsed the model. Rescission for fraud or mutual mistake does not do that. It instead reveals the boundary between repairing a licensed transition and refusing one that should not have been licensed.

\section{A tentative why}
This final thought is explicitly speculative. Contract law may have developed this structure because markets require both commitment and corrigibility. If the admissibility gate is too loose, obligation becomes noisy and coercive. If it is too strict, ordinary commerce cannot stabilize itself. If remedies are too thin, breach becomes a permission to externalize loss; if too rigid, every deviation becomes destructive. The law therefore evolves not around validity alone but around a managed balance between entry into obligation and legally intelligible continuation after failure.

That is why contract law is such a good companion to \textit{Contradiction and Continuation}. The field has, in ordinary doctrinal language, long operated with the practical insight that a state of affairs does not contain enough structure to explain its own continuation. Offer and acceptance start the relation, but breach and remedy tell us what the relation was capable of being all along.

\section{Conclusion}
Offer, acceptance, and consideration remain indispensable. But they are best understood not as isolated metaphysical ingredients of contracthood, but as admissibility conditions for entering a legally continuable relation. Material breach then appears as an inadmissible attempted continuation, substantial performance as a threshold doctrine for tolerable variation, and remedies as distinct repair operators preserving different invariants. Expectation damages preserve the promised future; reliance and restitution preserve narrower ones; rescission marks the point at which the law prefers unwinding to continuation. Read this way, classical contract doctrine is not displaced by the admissibility framework. It is rendered more structurally intelligible by it.

\end{document}
